{
\footnotesize
\setlength{\tabcolsep}{3pt}
\begin{longtable}{|>{\centering\arraybackslash}m{1.4cm}|>{\centering\arraybackslash}m{1.75cm}|m{5.9cm}|m{4cm}|}
\caption{Skenario Uji Penerimaan per Tujuan Platform}
\label{tab:uji_penerimaan} \\
\hline
\textbf{Skenario} & \textbf{Kebutuhan} & \multicolumn{1}{>{\centering\arraybackslash}m{5.9cm}|}{\textbf{Prosedur Uji}} & \multicolumn{1}{>{\centering\arraybackslash}m{4cm}|}{\textbf{Kriteria Penerimaan}} \\
\hline
\endfirsthead
\hline
\textbf{Skenario} & \textbf{Kebutuhan} & \multicolumn{1}{>{\centering\arraybackslash}m{5.9cm}|}{\textbf{Prosedur Uji}} & \multicolumn{1}{>{\centering\arraybackslash}m{4cm}|}{\textbf{Kriteria Penerimaan}} \\
\hline
\endhead
\endfoot
\hline
\endlastfoot

\multicolumn{4}{|l|}{\textbf{Tujuan T-1: arsitektur platform terintegrasi dan praktik GitOps}} \\
\hline
SK-F-01 & KF-01 & Mengubah definisi \textit{feature view} pada repositori Gitea lalu mengamati rekonsiliasi Argo CD, kemudian memverifikasi \texttt{feast feature-views list}. & \textit{Feature view} baru aktif tanpa intervensi manual dan perubahan terlacak pada riwayat Git. \\
\hline
SK-F-06 & KF-06 & Mempromosikan model melalui Argo Rollouts dengan strategi kanari, lalu menyuntikkan model bermasalah untuk memicu \textit{rollback}. & Model valid dipromosikan dan model bermasalah memicu \textit{rollback} otomatis tanpa kehilangan permintaan. \\
\hline
SK-F-12 & KF-12 & Mengubah jumlah \textit{replica} pada manifes YAML di Gitea, lalu mengamati rekonsiliasi Argo CD. & Status \textit{cluster} konvergen ke definisi Git dan tidak menyisakan \textit{drift}. \\
\hline
SK-F-13 & KF-13 & Menjalankan notebook satu \textit{environment} uv yang memuat SDK Feast, MLflow, dan klien \textit{inference} KServe v2, lalu memanggil ketiganya sebagai operasi baca independen. & Ketiga klien termuat dan berjalan dari satu \textit{environment} uv tanpa konfigurasi tambahan per layanan. \\
\hline
SK-F-14 & KF-14 & Mengirim sejumlah \textit{job} pelatihan melebihi kuota \textit{ClusterQueue} Kueue, lalu mengamati status admisi. & \textit{Job} di luar anggaran kuota berstatus \textit{Pending}. Admisi mengikuti \textit{nominalQuota} dan \textit{borrowingLimit}. \\
\hline
SK-N-03 & KNF-03 & Menaikkan beban dengan \texttt{load/hpa-keda-rampup.js} sambil mengamati HPA, KEDA \textit{ScaledObject}, dan jumlah \textit{pod}. & Jumlah \textit{replica} naik mengikuti beban dan \textit{throughput} bertambah mendekati linear. \\
\hline
SK-N-04 & KNF-04 & Menyuntikkan \textit{fault} dengan Chaos Mesh (\texttt{pod-chaos} dan \texttt{network-chaos}) lalu mengamati pemulihan. & Beban kerja pulih otomatis dengan RTO dan RPO di bawah ambang yang ditetapkan. \\
\hline
SK-N-08 & KNF-08 & Mengganti satu komponen (misalnya \textit{runtime} penyajian) melalui antarmuka CRD dan \textit{overlay} Kustomize. & Penggantian terisolasi pada antarmuka tanpa mengubah komponen lain. \\
\hline
SK-N-09 & KNF-09 & Mengakses seluruh konsol platform melalui satu sesi identitas dex serta memeriksa ketersediaan dokumentasi operasional pada repositori Gitea. & Seluruh konsol dapat dicapai melalui satu jalur autentikasi tanpa login ulang dan dokumentasi prosedur inti tersedia. \\
\hline
SK-N-10 & KNF-10 & Mengukur rasio kompresi ClickHouse, utilisasi sumber daya, dan atribusi biaya \textit{workload} pada OpenCost. & Rasio kompresi dan utilisasi memenuhi target serta biaya per \textit{workload} dapat diatribusikan. \\
\hline
SK-N-11 & KNF-11 & Menerapkan ulang seluruh manifes GitOps pada \textit{node} k3s baru dari satu sumber Git. & Seluruh komponen terekonstruksi tanpa intervensi manual selain \textit{bootstrap} Argo CD. \\
\hline
SK-N-12 & KNF-12 & Menjalankan \texttt{make usecase-crypto-test} untuk cakupan uji dan \texttt{make phase-full} pada \textit{cluster} bersih sambil mengukur waktu konvergensi. & Cakupan uji melewati ambang yang ditetapkan dan \texttt{phase-full} konvergen tanpa langkah manual. \\
\hline

\multicolumn{4}{|l|}{\textbf{Tujuan T-2: sub-sistem tata kelola data}} \\
\hline
SK-F-08 & KF-08 & Menelusuri \textit{lineage} pada DataHub dari topik sumber hingga keluaran prediksi melalui antarmuka katalog. & Rantai \textit{lineage} utuh dari sumber data sampai prediksi tersaji pada satu graf metadata. \\
\hline
SK-F-09 & KF-09 & Menjalankan dua \textit{run} pelatihan lalu mencatatnya pada MLflow, kemudian memutar ulang \textit{run} dari \textit{snapshot} data dan parameter tercatat. & \textit{Run} tercatat lengkap dan pemutaran ulang menghasilkan selisih metrik di bawah satu persen. \\
\hline
SK-F-10 & KF-10 & Memindahkan model melalui status \textit{none}, \textit{staging}, \textit{production}, dan \textit{archived} pada registri MLflow. & Transisi mengikuti aturan yang ditetapkan dan \textit{audit trail} terekam lengkap. \\
\hline
SK-N-05 & KNF-05 & Menjalankan \textit{checkpoint} Great Expectations pada \textit{batch} data, lalu menyuntikkan pelanggaran skema, rentang nilai, dan kelengkapan untuk menguji penegakan kontrak. & Setiap pelanggaran kontrak kualitas terdeteksi dan \textit{batch} cacat ditolak sebelum masuk lapis hilir. \\
\hline
SK-N-06 & KNF-06 & Menelusuri satu \textit{trace} OpenTelemetry pada Grafana Tempo dari ingestasi hingga prediksi. & Satu \textit{trace} menautkan lintasan kritis dan metrik, log, serta \textit{trace} dapat dirujuk silang. \\
\hline
SK-N-07 & KNF-07 & Memindai \textit{cluster} dengan Trivy dan memeriksa konfigurasi TLS pada \textit{endpoint} layanan. & TLS 1.3 \textit{in transit} dan AES-256 \textit{at rest} terverifikasi tanpa kerentanan kritis terbuka. \\
\hline

\multicolumn{4}{|l|}{\textbf{Tujuan T-3: deteksi \textit{drift} dan pelatihan ulang}} \\
\hline
SK-F-03 & KF-03 & Mengaudit \texttt{get\_historical\_features} pada \textit{dataset} pelatihan untuk memastikan jendela \textit{point-in-time} dipatuhi. & Tidak ada baris yang memakai informasi setelah \textit{event timestamp} acuan. \\
\hline
SK-F-05 & KF-05 & Menjalankan \textit{pipeline} pelatihan pada Kubeflow Pipelines yang membaca fitur dari tabel \textit{gold} dan menulis model ke MLflow. & \textit{Pipeline} berjalan dari ujung ke ujung dan model tercatat sebagai artefak \textit{run} pada MLflow yang dapat disajikan KServe. \\
\hline
SK-F-07 & KF-07 & Menyuntikkan \textit{drift} terkendali pada satu fitur data nyata, lalu mengamati detektor \textit{drift} dan \textit{CronWorkflow retrain-on-drift}. & \textit{Drift} terdeteksi saat ambang PSI atau KS terlampaui dan \textit{pipeline} pelatihan ulang terpicu otomatis. \\
\hline

\multicolumn{4}{|l|}{\textbf{Tujuan T-4: layanan fitur \textit{dual-store}}} \\
\hline
SK-F-02 & KF-02 & Memeriksa penyajian fitur \textit{online} pada Valkey mengembalikan nilai non-\textit{null} melalui \texttt{online-serving-check.py}. & Penyajian \textit{online} mengembalikan nilai non-\textit{null} dan kesetaraan nilai \textit{online} dan \textit{offline} terpenuhi secara struktural melalui \textit{materialize}. \\
\hline
SK-F-04 & KF-04 & Menjalankan beban pencarian vektor pada Qdrant melalui \texttt{load/vector-search-latency.js} dengan k6. & Latensi pencarian vektor berada pada orde sub-detik; pengujian \textit{recall} dilakukan saat koleksi vektor terisi. \\
\hline
SK-F-11 & KF-11 & Membandingkan fitur hasil \textit{batch} (dbt) dan hasil \textit{stream} (Flink) untuk indikator yang sama. & Nilai fitur \textit{batch} dan \textit{stream} setara dalam toleransi yang ditetapkan. \\
\hline
SK-N-01 & KNF-01 & Menjalankan \texttt{load/feast-online-latency.js} dan beban pencarian vektor dengan k6 sambil memantau SLO Sloth. & Latensi p99 berada pada orde sub-detik pada \textit{feature serving} dan \textit{vector search}. \\
\hline
SK-N-02 & KNF-02 & Menjalankan \texttt{load/feast-throughput.js} untuk \textit{serving} dan \texttt{kafka-producer-perf-test} untuk \textit{stream}, sambil memantau \textit{lag}. & \textit{Throughput} mencapai target dan \textit{lag} konsumer tetap terkendali pada beban puncak. \\
\hline

\end{longtable}
}
